% !TeX program = lualatex
% Definition and sources only; this file is not an LLM solution attempt.
\documentclass[11pt,a4paper]{article}
\usepackage[margin=25mm]{geometry}
\usepackage{fontspec}
\setmainfont{lmroman10-regular.otf}[BoldFont=lmroman10-bold.otf,
 ItalicFont=lmroman10-italic.otf,BoldItalicFont=lmroman10-bolditalic.otf]
\usepackage{amsmath,amssymb,mathtools}
\usepackage{unicode-math}
\setmathfont{latinmodern-math.otf}
\AtBeginDocument{\let\setminus\smallsetminus}
\usepackage{xurl,hyperref}
\hypersetup{colorlinks=true,urlcolor=blue}
\setcounter{secnumdepth}{0}
\setlength{\parindent}{0pt}
\setlength{\parskip}{5pt}
\setlength{\emergencystretch}{3em}
\begin{document}
\section*{Questions in Geometric Group Theory — Q 12.7}\label{statement}

\subsection{Definitions and mathematical statement}
Let $F_n$ be the free group on a fixed basis of size $n\geq2$, and let $\alpha\in\operatorname{Aut}(F_n)$. A right-invariant order is a strict total order $<$ on $F_n$ such that
\[
 g<h\ \Longrightarrow\ gk<hk\qquad(g,h,k\in F_n).
\]
The automorphism preserves this order when $g<h$ implies $\alpha(g)<\alpha(h)$. Characterize the automorphisms for which some right-invariant order is preserved.

Equivalently, seek a subset $P\subset F_n$ satisfying
\[
 F_n=P\sqcup\{1\}\sqcup P^{-1},\qquad PP\subset P,
 \qquad\alpha(P)=P.
\]
The corresponding order is $g<h$ exactly when $hg^{-1}\in P$. No invariance of $P$ under conjugation is imposed, since that would require a two-sided order and would change the question.

A necessary obstruction is a nonconstant finite orbit: if $\alpha^m(g)=g$ for some $m>1$ but $\alpha(g)\ne g$, order preservation would give an impossible strict cycle of inequalities. Fixed elements themselves are allowed; in particular the identity automorphism preserves every right order. Determine whether further obstructions are needed and, if so, formulate a useful necessary-and-sufficient condition from the automorphism. The question concerns an order on the entire free group, not just an order on its abelianization. The chosen order may depend on the automorphism and on its rank.
\subsection{Short English statement}
Which automorphisms of a free group preserve a right-invariant total order?
\subsection{Sources}
\begin{itemize}
\item[S1] ULAM.AI and the UnsolvedMath Contributors, supplied problems.json, catalogue entry, AMR Open Problem Lists. Catalogue identity and source transcription; CC BY 4.0.
\url{https://huggingface.co/datasets/ulamai/UnsolvedMath}.
\item[S2] Bestvina - Questions in Geometric Group Theory (pdf) (2004), Question 12.7 (PDF page 20). 
\url{https://www.math.utah.edu/~bestvina/eprints/questions-updated.pdf}.
\end{itemize}
\end{document}
